\documentclass[12pt, a4paper]{article}

% ── Codifica e lingua ───────────────────────────────────────────────────────
\usepackage[utf8]{inputenc}
\usepackage[T1]{fontenc}
\usepackage[italian]{babel}

% ── Matematica ──────────────────────────────────────────────────────────────
\usepackage{amsmath, amssymb, amsthm}
\usepackage{mathtools}          % dcases, prescript, ecc.
\usepackage{physics}            % \bra, \ket, \norm, \abs, \dv, \pdv

% ── Figure e layout ─────────────────────────────────────────────────────────
\usepackage{graphicx}
\usepackage{float}
\usepackage{subcaption}
\usepackage[margin=2.5cm]{geometry}
\usepackage{setspace}
\onehalfspacing

% ── Hyperlink ───────────────────────────────────────────────────────────────
\usepackage[colorlinks=true, linkcolor=blue!70!black,
            citecolor=blue!70!black, urlcolor=blue!70!black]{hyperref}

% ── Listati di codice ────────────────────────────────────────────────────────
\usepackage{listings}
\usepackage{xcolor}
\lstset{
  basicstyle=\ttfamily\small,
  keywordstyle=\color{blue!80!black}\bfseries,
  commentstyle=\color{gray},
  stringstyle=\color{green!50!black},
  breaklines=true,
  frame=single,
  numbers=left,
  numberstyle=\tiny\color{gray}
}

% ── Comandi personalizzati ──────────────────────────────────────────────────
\newcommand{\dt}{\Delta t}
\newcommand{\dx}{\Delta x}
\newcommand{\hb}{\hbar}
\newcommand{\ii}{\mathrm{i}}          % unità immaginaria (tonda)
\newcommand{\ee}{\mathrm{e}}          % numero di Eulero (tondo)
\newcommand{\CN}{Crank--Nicolson}
\newcommand{\Htop}{\hat{H}}
\newcommand{\Ktop}{\hat{K}}
\newcommand{\Vtop}{\hat{V}}

% ── Bibliografia ─────────────────────────────────────────────────────────────
\usepackage[numbers, sort&compress]{natbib}

% ─────────────────────────────────────────────────────────────────────────────
\begin{document}

% ── Frontespizio ─────────────────────────────────────────────────────────────
\begin{titlepage}
  \centering
  \vspace*{3cm}
  {\LARGE\bfseries Soluzione Numerica dell'Equazione di Schr\"{o}dinger\\[6pt]
   Dipendente dal Tempo\\[4pt]
   tramite lo Schema di Crank--Nicolson\par}
  \vspace{1.5cm}
  {\large Effetto tunnel quantistico di un pacchetto d'onde gaussiano\\
   attraverso una barriera di potenziale rettangolare\par}
  \vspace{2cm}
  {\large Ludovico Fecia di Cossato\par}
  \vfill
  {\large\itshape TDSE\_CN — note teoriche\par}
  \vspace{0.5cm}
  {\normalsize Unità atomiche ovunque: $\hbar = m_e = e = 1$\par}
  \vspace{1cm}
  {\normalsize \today\par}
\end{titlepage}

\tableofcontents
\newpage

% ═══════════════════════════════════════════════════════════════════════════
\section{Introduzione}
% ═══════════════════════════════════════════════════════════════════════════

L'equazione di Schr\"{o}dinger dipendente dal tempo (TDSE) governa l'evoluzione
di qualsiasi sistema quantistico non relativistico. Anche in una sola dimensione
spaziale, essa presenta fenomeni — allargamento del pacchetto d'onde, interferenza
ed effetto tunnel quantistico — che non hanno analoghi classici e sono centrali
nella fisica moderna e nella tecnologia (diodi tunnel, microscopia a
scansione effetto tunnel, decadimento alfa nucleare, \ldots).

Nonostante la sua profondità concettuale, la TDSE unidimensionale con potenziale
statico è completamente trattabile numericamente: il problema si riduce a propagare
un vettore complesso nel tempo tramite una sequenza di sistemi lineari sparsi.
Questo la rende un banco di prova ideale per i metodi alle differenze finite
impliciti e per sviluppare l'intuizione sulla dinamica quantistica.

Queste note descrivono il problema fisico, la struttura matematica della
discretizzazione \CN{} e l'algoritmo di Thomas utilizzato per risolvere il
sistema tridiagonale risultante ad ogni passo temporale. Una sezione finale
analizza l'output della simulazione numerica in modo qualitativo e quantitativo.

% ═══════════════════════════════════════════════════════════════════════════
\section{Impostazione fisica}
% ═══════════════════════════════════════════════════════════════════════════

\subsection{L'equazione di Schr\"{o}dinger dipendente dal tempo}

Lo stato di una particella quantistica unidimensionale è completamente descritto
da una funzione d'onda a valori complessi $\psi(x,t) \in L^2(\mathbb{R})$.
La sua evoluzione temporale è governata dalla TDSE
\begin{equation}
  \ii\hbar\,\pdv{\psi}{t}(x,t) \;=\; \Htop\,\psi(x,t),
  \label{eq:tdse}
\end{equation}
dove l'operatore hamiltoniano $\Htop = \Ktop + \Vtop$ si divide in una
parte cinetica
\begin{equation}
  \Ktop \;=\; -\frac{\hbar^2}{2m}\pdv[2]{}{x}
\end{equation}
e un operatore di energia potenziale $\Vtop = V(x)$ che agisce in modo
moltiplicativo.

In \textbf{unità atomiche} ($\hbar = m_e = 1$) — usate in tutto questo
documento — l'equazione \eqref{eq:tdse} diventa
\begin{equation}
  \ii\,\pdv{\psi}{t} \;=\;
  \left[-\frac{1}{2}\pdv[2]{}{x} + V(x)\right]\psi.
  \label{eq:tdse_au}
\end{equation}

\subsection{Interpretazione fisica della funzione d'onda}

La funzione d'onda $\psi(x,t)$ non è direttamente osservabile. La quantità
fisicamente significativa è la \emph{densità di probabilità}
\begin{equation}
  \rho(x,t) \;=\; |\psi(x,t)|^2,
\end{equation}
che fornisce la probabilità di trovare la particella tra $x$ e $x+dx$
come $\rho(x,t)\,dx$. La probabilità totale è conservata:
\begin{equation}
  \int_{-\infty}^{+\infty} |\psi(x,t)|^2\,dx \;=\; 1
  \quad \forall\, t.
  \label{eq:norm}
\end{equation}
Questo vincolo — la \emph{conservazione della norma} o \emph{unitarietà}
dell'evoluzione — è uno dei controlli di consistenza più importanti per
qualsiasi schema numerico.

\subsection{Il pacchetto d'onde gaussiano}

Lo stato iniziale scelto per questa simulazione è un \emph{pacchetto d'onde
gaussiano}, lo stato quantistico che più si avvicina a una particella classica
con posizione e momento definiti:
\begin{equation}
  \psi(x,0) \;=\; A\,\exp\!\left[-\frac{(x-x_0)^2}{4\sigma^2}\right]
  \exp\!\left(\ii k_0 x\right),
  \label{eq:gauss}
\end{equation}
dove:
\begin{itemize}
  \item $x_0$ è la posizione iniziale del centro del pacchetto,
  \item $\sigma$ è la larghezza spaziale (deviazione standard di $\rho$),
  \item $k_0$ è il vettore d'onda centrale, corrispondente al momento medio
        $\langle p\rangle = \hbar k_0$,
  \item $A = (2\pi\sigma^2)^{-1/4}$ è la costante di normalizzazione.
\end{itemize}

Il gaussiano è lo \emph{stato di minima indeterminazione}: esso satura
la relazione di indeterminazione di Heisenberg $\Delta x\,\Delta p \geq \hbar/2$.
La sua trasformata di Fourier è anch'essa un gaussiano, centrato in $k_0$
con larghezza $1/(2\sigma)$, così che la dispersione del momento è
$\Delta p = \hbar/(2\sigma)$.

L'energia cinetica media del pacchetto è
\begin{equation}
  E \;=\; \frac{\langle p\rangle^2}{2m} \;=\; \frac{k_0^2}{2}
  \quad \text{(unità atomiche)}.
  \label{eq:energy}
\end{equation}
Questa è la quantità da confrontare con l'altezza della barriera $V_0$ per
classificare il regime dinamico.

\subsection{La barriera di potenziale rettangolare}

Il potenziale esterno è una barriera rettangolare di altezza $V_0$ e
larghezza $w$ situata in $x \in [x_\text{bar},\, x_\text{bar}+w]$:
\begin{equation}
  V(x) \;=\;
  \begin{cases}
    V_0 & \text{se } x_\text{bar} < x < x_\text{bar} + w, \\
    0   & \text{altrimenti.}
  \end{cases}
  \label{eq:barrier}
\end{equation}
La barriera è il potenziale di scattering non banale più semplice e ammette
una soluzione analitica esatta per stati a onda piana (si veda la
Sezione~\ref{sec:tunnelling}), rendendola ideale per validare il metodo
numerico.

% ═══════════════════════════════════════════════════════════════════════════
\section{Effetto tunnel quantistico}
\label{sec:tunnelling}
% ═══════════════════════════════════════════════════════════════════════════

\subsection{Analisi stazionaria}

Per un'onda piana monocromatica $\psi(x) = \ee^{\ii kx}$ con energia
$E = k^2/2$, l'equazione di Schr\"{o}dinger può essere risolta a tratti
in ciascuna regione. Il raccordo delle funzioni d'onda e delle loro
derivate ai bordi della barriera fornisce il coefficiente di trasmissione
esatto \cite{Griffiths2018}
\begin{equation}
  T \;=\; \left[1 + \frac{V_0^2\sinh^2(\kappa w)}{4E(V_0-E)}\right]^{-1}
  \qquad (E < V_0),
  \label{eq:T_tunnelling}
\end{equation}
dove $\kappa = \sqrt{2(V_0-E)}$ (unità atomiche) è la costante di decadimento
evanescente all'interno della barriera. Il coefficiente di riflessione è
semplicemente $R = 1 - T$.

Dall'equazione \eqref{eq:T_tunnelling} emergono immediatamente alcune
caratteristiche qualitative:
\begin{enumerate}
  \item $T > 0$ anche quando $E < V_0$ — questo è l'\emph{effetto tunnel quantistico}.
        Classicamente, la particella non può penetrare in una regione in cui $E < V$.
  \item $T$ decade \emph{esponenzialmente} con la larghezza della barriera:
        per $\kappa w \gg 1$,
        \begin{equation}
          T \;\approx\; 16\,\frac{E}{V_0}\!\left(1 - \frac{E}{V_0}\right)
          \ee^{-2\kappa w}.
        \end{equation}
  \item Nel regime sopra-barriera $E > V_0$, si verifica ancora una riflessione
        parziale dovuta alla natura ondulatoria di $\psi$, con
        \begin{equation}
          T \;=\; \left[1 + \frac{V_0^2\sin^2(k' w)}{4E(E-V_0)}\right]^{-1},
          \quad k' = \sqrt{2(E-V_0)}.
        \end{equation}
        La trasmissione risonante ($T = 1$) si verifica ogniqualvolta $k'w = n\pi$.
\end{enumerate}

\subsection{Dinamica del pacchetto d'onde}

Un pacchetto d'onde gaussiano è una sovrapposizione di onde piane:
\begin{equation}
  \psi(x,0) \;=\; \int_{-\infty}^{+\infty}
  \tilde{\phi}(k)\,\ee^{\ii kx}\,\frac{dk}{2\pi},
  \qquad
  \tilde{\phi}(k) \propto \ee^{-(k-k_0)^2\sigma^2}\,\ee^{-\ii k x_0}.
\end{equation}
Ciascuna componente $k$ evolve indipendentemente e interagisce con la barriera
con il proprio coefficiente $T(k)$. Dopo l'evento di scattering, la funzione
d'onda totale si divide in una parte \emph{riflessa} che viaggia verso sinistra
e una parte \emph{trasmessa} che viaggia verso destra.

Le probabilità integrate
\begin{align}
  R(t) &= \int_{-\infty}^{x_\text{bar}} |\psi(x,t)|^2\,dx, \\
  T(t) &= \int_{x_\text{bar}}^{+\infty} |\psi(x,t)|^2\,dx
\end{align}
si avvicinano a valori asintotici costanti al termine dello scattering.
Questi valori asintotici possono essere confrontati con la formula a
onda piana \eqref{eq:T_tunnelling} valutata in $k = k_0$ — si attende
un accordo quando il pacchetto è stretto nello spazio $k$ (grande $\sigma$).

% ═══════════════════════════════════════════════════════════════════════════
\section{Metodo numerico}
% ═══════════════════════════════════════════════════════════════════════════

\subsection{Discretizzazione spaziale}

Il dominio $[x_\text{min},\, x_\text{max}]$ è suddiviso in $N$ punti
uniformemente spaziati con passo $\dx = L/(N-1)$, dove $L =
x_\text{max} - x_\text{min}$. La funzione d'onda è rappresentata come
il vettore $\boldsymbol{\psi}^n \in \mathbb{C}^N$ con componenti
$\psi_i^n \approx \psi(x_i, n\dt)$.

La derivata seconda è sostituita dallo stencil standard alle differenze
finite simmetriche del secondo ordine:
\begin{equation}
  \pdv[2]{\psi}{x}\bigg|_{x_i}
  \;\approx\;
  \frac{\psi_{i-1} - 2\psi_i + \psi_{i+1}}{\dx^2}
  + \mathcal{O}(\dx^2).
  \label{eq:fd_stencil}
\end{equation}
Questo approssima la parte cinetica di $\Htop$ come una matrice tridiagonale
reale simmetrica $K$ con diagonale principale $-2/\dx^2$ e sottodiagonali
$+1/\dx^2$ (moltiplicato per $-1/2$):
\begin{equation}
  [H_\text{disc}]_{ij}
  \;=\;
  -\frac{1}{2\dx^2}\begin{pmatrix}
    -2 & 1 \\
    1 & -2 & 1 \\
      & \ddots & \ddots & \ddots \\
      &        & 1 & -2 & 1 \\
      &        &   & 1  & -2
  \end{pmatrix}_{ij}
  + V_i\,\delta_{ij}.
  \label{eq:Hdisc}
\end{equation}

\subsection{Discretizzazione temporale: lo schema di Crank--Nicolson}

Una discretizzazione esplicita ingenua (Eulero in avanti) della TDSE,
\begin{equation}
  \boldsymbol{\psi}^{n+1} = \boldsymbol{\psi}^n
  - \ii\dt\, H_\text{disc}\,\boldsymbol{\psi}^n,
\end{equation}
è \emph{condizionatamente stabile}: il passo temporale deve soddisfare
$\dt \lesssim \dx^2$, il che per $\dx = 0{,}01$\,bohr significa
$\dt \lesssim 10^{-4}$\,u.a. — due ordini di grandezza inferiori al
passo \CN{} usato qui. Inoltre, lo schema esplicito non conserva
la norma esattamente.

Lo \textbf{schema \CN{}} \cite{CrankNicolson1947} è un metodo implicito
del secondo ordine che media il membro destro tra il livello temporale
corrente e quello successivo:
\begin{equation}
  \ii\,\frac{\boldsymbol{\psi}^{n+1} - \boldsymbol{\psi}^n}{\dt}
  \;=\;
  \frac{H_\text{disc}}{2}\!\left(\boldsymbol{\psi}^{n+1}
  + \boldsymbol{\psi}^n\right).
  \label{eq:cn_raw}
\end{equation}
Riorganizzando \eqref{eq:cn_raw}:
\begin{equation}
  \underbrace{\left(I + \frac{\ii\dt}{2}\,H_\text{disc}\right)}_{A}
  \boldsymbol{\psi}^{n+1}
  \;=\;
  \underbrace{\left(I - \frac{\ii\dt}{2}\,H_\text{disc}\right)}_{B}
  \boldsymbol{\psi}^n,
  \label{eq:cn_system}
\end{equation}
che è un sistema lineare $A\,\boldsymbol{\psi}^{n+1} = \mathbf{b}^n$
con $\mathbf{b}^n = B\,\boldsymbol{\psi}^n$.

\subsection{Proprietà dello schema di Crank--Nicolson}
\label{sec:cn_properties}

\paragraph{Stabilità incondizionata.}
Gli autovalori di $H_\text{disc}$ sono reali (è hermitiano). Per un
autovalore $\lambda$, il fattore di amplificazione dello schema \CN{} è
\begin{equation}
  g \;=\; \frac{1 - \ii\dt\lambda/2}{1 + \ii\dt\lambda/2},
\end{equation}
che soddisfa $|g| = 1$ per \emph{qualsiasi} $\dt > 0$. Lo schema è
quindi \textbf{incondizionatamente stabile}: non è richiesta alcuna
restrizione sul passo temporale per la stabilità.

\paragraph{Conservazione della norma (probabilità).}
Le matrici $A$ e $B$ soddisfano $B = A^\dagger$ (trasposta coniugata
complessa). Il propagatore discreto $A^{-1}B$ è quindi \emph{unitario}:
\begin{equation}
  (A^{-1}B)^\dagger (A^{-1}B)
  \;=\; B^\dagger (A^\dagger)^{-1} A^{-1} B
  \;=\; A A^{-1} (A^\dagger)^{-1} B
  \;=\; I.
\end{equation}
Di conseguenza, $\|\boldsymbol{\psi}^{n+1}\|^2 = \|\boldsymbol{\psi}^n\|^2$
ad ogni passo — la norma discreta è conservata esattamente (a meno di
errori di arrotondamento in virgola mobile dell'ordine di $10^{-14}$).

\paragraph{Accuratezza del secondo ordine.}
L'errore di troncamento locale di \eqref{eq:cn_raw} è $\mathcal{O}(\dt^2)$
nel tempo e $\mathcal{O}(\dx^2)$ nello spazio, quindi lo schema è
globalmente del secondo ordine in entrambe le variabili.

\subsection{Struttura delle matrici}

Le matrici $A$ e $B$ sono \emph{tridiagonali complesse}. Definendo il
parametro
\begin{equation}
  r \;=\; \frac{\dt}{4\dx^2},
\end{equation}
i loro elementi non nulli sono (per $1 \leq i \leq N-2$):
\begin{alignat}{2}
  A_{i,i}   &= 1 + \ii\frac{\dt}{2}\!\left(\frac{1}{\dx^2} + V_i\right),
  &\quad A_{i,i\pm 1} &= -\ii r, \label{eq:A_elements}\\[4pt]
  B_{i,i}   &= 1 - \ii\frac{\dt}{2}\!\left(\frac{1}{\dx^2} + V_i\right),
  &\quad B_{i,i\pm 1} &= +\ii r. \label{eq:B_elements}
\end{alignat}
La prima e l'ultima riga impongono le condizioni al contorno di Dirichlet
$\psi_0 = \psi_{N-1} = 0$: entrambe le righe vengono impostate alla
riga identità $(1,\,0,\,\ldots)$ in $A$ e $B$.

Poiché $V(x)$ è statico, $A$ e $B$ vengono \textbf{costruite una sola volta}
prima del ciclo temporale e riutilizzate ad ogni passo.

\subsection{L'algoritmo di Thomas}
\label{sec:thomas}

Ad ogni passo temporale, il membro destro viene calcolato come prodotto
matrice-vettore $\mathbf{b}^n = B\,\boldsymbol{\psi}^n$ (una passata sulla
griglia, $\mathcal{O}(N)$ operazioni), e il sistema tridiagonale risultante
$A\,\boldsymbol{\psi}^{n+1} = \mathbf{b}^n$ viene risolto con
l'\textbf{algoritmo di Thomas} \cite{Thomas1949}.

L'algoritmo di Thomas è l'eliminazione gaussiana specializzata per sistemi
tridiagonali. Per un sistema con sottodiagonale $l_i$, diagonale principale
$d_i$, sopradiagonale $u_i$ e membro destro $b_i$:

\paragraph{Passata in avanti} (eliminazione di $l_i$, da $i=1$ a $N-1$):
\begin{align}
  w_i &= l_i / d_{i-1}, \\
  d_i &\leftarrow d_i - w_i\,u_{i-1}, \\
  b_i &\leftarrow b_i - w_i\,b_{i-1}.
\end{align}

\paragraph{Sostituzione all'indietro} (da $i=N-1$ fino a $0$):
\begin{align}
  x_{N-1} &= b_{N-1}/d_{N-1}, \\
  x_i     &= (b_i - u_i\,x_{i+1}) / d_i.
\end{align}

Il costo totale è $\mathcal{O}(N)$ sia in tempo che in memoria — ottimale per
un sistema tridiagonale. L'algoritmo è numericamente stabile per matrici
a diagonale dominante, condizione che la matrice \CN{} $A$ soddisfa sempre
(la diagonale principale domina perché $|A_{ii}| > 2|A_{i,i\pm 1}|$
ogniqualvolta $\dt/\dx^2 < 2$, condizione facilmente verificata in pratica).

Questo è lo stesso algoritmo implementato dalla routine LAPACK \texttt{zgtsv},
qui realizzato esplicitamente in aritmetica complessa senza alcuna dipendenza
da librerie esterne.

\subsection{Condizioni al contorno}

Le condizioni al contorno di \textbf{Dirichlet} $\psi(x_\text{min},t) =
\psi(x_\text{max},t) = 0$ sono imposte. Fisicamente, le pareti agiscono
come barriere di potenziale infinite. Introducono riflessioni artificiali
nel momento in cui il pacchetto d'onde raggiunge il bordo, quindi la
lunghezza del dominio $L$ deve essere scelta sufficientemente grande
affinché il pacchetto rimanga lontano dalle pareti per tutta la simulazione.

Un criterio pratico è $x_\text{min} \ll x_0 - k_0 T_\text{tot}$ e
$x_\text{max} \gg x_0 + k_0 T_\text{tot}$, dove $k_0 T_\text{tot}$ è
la distanza approssimativa percorsa dal centro del pacchetto.

\subsection{Schema dell'algoritmo}

L'algoritmo completo di avanzamento temporale è:

\begin{enumerate}
  \item Leggere i parametri e costruire la griglia $\{x_i\}$ e il potenziale $\{V_i\}$.
  \item Inizializzare $\boldsymbol{\psi}^0$ come il gaussiano normalizzato
        \eqref{eq:gauss}.
  \item Costruire le matrici tridiagonali $A$ e $B$ (una sola volta).
  \item \textbf{Per} $n = 0, 1, \ldots, M-1$:
  \begin{enumerate}
    \item[(a)] Calcolare $\mathbf{b}^n = B\,\boldsymbol{\psi}^n$
               (prodotto matrice-vettore, $\mathcal{O}(N)$).
    \item[(b)] Risolvere $A\,\boldsymbol{\psi}^{n+1} = \mathbf{b}^n$
               con l'algoritmo di Thomas ($\mathcal{O}(N)$).
    \item[(c)] Imporre $\psi_0^{n+1} = \psi_{N-1}^{n+1} = 0$.
    \item[(d)] Scrivere istantanea e norma su file (se è un passo di istantanea).
  \end{enumerate}
\end{enumerate}
Costo totale per passo temporale: $\mathcal{O}(N)$. Costo totale: $\mathcal{O}(MN)$.

% ═══════════════════════════════════════════════════════════════════════════
\section{Analisi dell'output della simulazione}
% ═══════════════════════════════════════════════════════════════════════════

Questa sezione descrive il comportamento fisico atteso e le corrispondenti
signature nei tre grafici prodotti dagli script Python di analisi.

I parametri di default della simulazione sono:
\begin{center}
\begin{tabular}{lll}
\hline
Parametro & Valore & Ruolo \\
\hline
$L$ & 20\,bohr & lunghezza del dominio \\
$\dx$ & 0,01\,bohr & passo spaziale \\
$\dt$ & 0,001\,u.a. & passo temporale \\
$T_\text{tot}$ & 2,0\,u.a. & tempo totale \\
$x_0$ & $-5{,}0$\,bohr & centro iniziale del pacchetto \\
$\sigma$ & 0,5\,bohr & larghezza del pacchetto \\
$k_0$ & 5,0\,bohr$^{-1}$ & vettore d'onda centrale \\
$x_\text{bar}$ & 0,0\,bohr & bordo sinistro della barriera \\
$w$ & 0,5\,bohr & larghezza della barriera \\
$V_0$ & 20,0\,$E_h$ & altezza della barriera \\
\hline
\end{tabular}
\end{center}

L'energia cinetica media del pacchetto è $E = k_0^2/2 = 12{,}5\,E_h$, mentre
l'altezza della barriera è $V_0 = 20{,}0\,E_h$, quindi $E < V_0$: la simulazione
è nel \textbf{regime di tunnel}. La costante di decadimento evanescente
all'interno della barriera è $\kappa = \sqrt{2(V_0-E)} \approx 3{,}87$\,bohr$^{-1}$,
fornendo una probabilità di trasmissione approssimativa dall'equazione
\eqref{eq:T_tunnelling}:
\begin{equation}
  T_\text{analitico} \;\approx\;
  \left[1 + \frac{20^2\sinh^2(3{,}87\times 0{,}5)}{4\times 12{,}5\times 7{,}5}
  \right]^{-1}
  \;\approx\; 0{,}04.
  \label{eq:T_predicted}
\end{equation}
Si prevede che circa il 4\% della probabilità effettui il tunnel attraverso
la barriera, con il restante 96\% riflesso.

\subsection{Grafico delle istantanee (\texttt{snapshots.png})}

\paragraph{Cosa viene mostrato.}
Cinque pannelli visualizzano $|\psi(x,t)|^2$ a tempi equidistanti
$t = 0,\, T/4,\, T/2,\, 3T/4,\, T$.
La barriera di potenziale è mostrata come una regione ombreggiata scalata
alla densità di probabilità massima come riferimento visivo.

\paragraph{Comportamento atteso.}
\begin{itemize}
  \item \textbf{$t = 0$}: Un singolo picco gaussiano centrato in $x_0 = -5$\,bohr,
        ben a sinistra della barriera. La sua larghezza è $\sigma = 0{,}5$\,bohr.

  \item \textbf{Fase di propagazione ($t \approx T/4$)}: Il pacchetto si muove
        verso destra con velocità di gruppo $v_g = k_0 = 5$\,bohr/u.a. Si
        allarga gradualmente per dispersione: poiché componenti $k$ diverse
        viaggiano a velocità leggermente diverse $v(k) = k$, il pacchetto
        si allarga come $\sigma(t) \approx \sqrt{\sigma^2 + t^2/(4\sigma^2)}$.

  \item \textbf{Scattering ($t \approx T/2$)}: Il pacchetto raggiunge la
        barriera. La densità di probabilità viene temporaneamente compressa
        e distorta — il bordo di attacco decelera mentre la coda recupera.
        Una piccola frazione trasmessa filtra attraverso il lato destro
        della barriera.

  \item \textbf{Post-scattering ($t \approx 3T/4$ e $T$)}: Due picchi distinti
        sono visibili. Un grande pacchetto riflesso viaggia verso sinistra;
        un piccolo pacchetto trasmesso continua verso destra. Le aree relative
        sotto i due picchi riflettono il rapporto $R:T \approx 96:4$.
\end{itemize}

\paragraph{Cosa osservare.}
L'area totale sotto la curva (norma) deve rimanere costante in tutti i
pannelli. Qualsiasi crescita o decadimento visibile segnala un problema
numerico.

\subsection{Grafico dell'analisi quantitativa (\texttt{analysis.png})}

\subsubsection{Pannello A: Conservazione della norma}

\paragraph{Cosa viene mostrato.}
La norma discreta $\|\boldsymbol{\psi}^n\|^2_\dx = \dx\sum_i |\psi_i^n|^2$
in funzione del tempo, confrontata con il valore esatto di~1.

\paragraph{Comportamento atteso.}
Poiché il propagatore \CN{} è unitario (Sezione~\ref{sec:cn_properties}),
la deviazione della norma dovrebbe essere $|\|\psi^n\|^2 - 1| \lesssim 10^{-12}$
durante tutta la simulazione — essenzialmente la precisione della macchina.
Questo è fondamentalmente diverso dagli schemi espliciti, che dissipano o
amplificano la norma.

\paragraph{Cosa osservare.}
Una linea piatta a 1 conferma l'unitarietà del propagatore numerico.
Qualsiasi deriva monotona indicherebbe un bug nell'implementazione; una
deriva oscillatoria suggerisce problemi di stabilità.

\subsubsection{Pannello B: Coefficienti di trasmissione e riflessione}

\paragraph{Cosa viene mostrato.}
L'evoluzione temporale di
\begin{equation}
  R(t) = \dx\sum_{i:\,x_i < x_\text{bar}} |\psi_i|^2,
  \qquad
  T(t) = \dx\sum_{i:\,x_i \geq x_\text{bar}} |\psi_i|^2,
\end{equation}
insieme alla loro somma $R(t) + T(t)$.

\paragraph{Comportamento atteso.}
Inizialmente $R \approx 1$ e $T \approx 0$ (l'intero pacchetto a sinistra
della barriera). Quando il pacchetto raggiunge e interagisce con la barriera,
$R$ diminuisce e $T$ aumenta. Al termine dello scattering, entrambi i
coefficienti raggiungono i loro valori asintotici: $T_\infty \approx 0{,}04$
e $R_\infty \approx 0{,}96$ da \eqref{eq:T_predicted}.

La somma $R(t) + T(t)$ deve rimanere uguale a 1 in ogni istante (è solo
un altro modo di scrivere la conservazione della norma), fornendo un
controllo di autoconsistenza. Durante la fase di scattering, $R+T$ può
apparire leggermente ridotta se la funzione d'onda è distribuita
\emph{attraverso} il bordo della barriera, ma questo è un effetto di
bordo della regione di integrazione, non una violazione fisica.

\paragraph{Cosa osservare.}
Il valore asintotico di $T_\infty$ può essere direttamente confrontato
con il risultato analitico a onda piana \eqref{eq:T_predicted}. Le discrepanze
sorgono perché il pacchetto gaussiano ha una dispersione di momento finita;
il risultato del pacchetto d'onde converge alla formula a onda piana nel
limite $\sigma \to \infty$ (distribuzione stretta in $k$).

\subsubsection{Pannello C: Posizione media $\langle x \rangle(t)$}

\paragraph{Cosa viene mostrato.}
Il valore di aspettazione della posizione,
\begin{equation}
  \langle x\rangle(t) \;=\; \dx\sum_i x_i\,|\psi_i(t)|^2,
\end{equation}
in funzione del tempo.

\paragraph{Comportamento atteso.}
\begin{itemize}
  \item \textbf{Prima della barriera}: $\langle x\rangle$ aumenta linearmente
        con pendenza $v_g = k_0 = 5$\,bohr/u.a., riflettendo il moto
        classico del centro del pacchetto.

  \item \textbf{Durante lo scattering}: la pendenza diminuisce mentre il
        pacchetto interagisce con la barriera. Classicamente, la particella
        sarebbe riflessa istantaneamente; quantisticamente, l'interazione
        è distribuita su un tempo finito (relativo al \emph{tempo di ritardo
        di Wigner}).

  \item \textbf{Dopo lo scattering}: $\langle x\rangle$ è la media pesata
        per probabilità dei centri riflesso e trasmesso. Poiché $R \gg T$,
        la posizione media si sposta verso sinistra (il pacchetto riflesso
        domina). La pendenza del regime lineare asintotico sarà negativa,
        approssimativamente $-k_0 R_\infty$.
\end{itemize}

\paragraph{Cosa osservare.}
La transizione da pendenza positiva a negativa segnala il momento dello
scattering. La pendenza asintotica codifica il peso relativo delle
componenti riflessa e trasmessa.

\subsubsection{Pannello D: Larghezza del pacchetto $\sigma(t)$}

\paragraph{Cosa viene mostrato.}
La larghezza quadratica media del pacchetto,
\begin{equation}
  \sigma(t) \;=\; \sqrt{\langle x^2\rangle - \langle x\rangle^2},
  \quad
  \langle x^2\rangle = \dx\sum_i x_i^2\,|\psi_i|^2.
\end{equation}

\paragraph{Comportamento atteso.}
\begin{itemize}
  \item \textbf{Propagazione libera}: Un pacchetto d'onde gaussiano nello
        spazio libero si allarga come
        \begin{equation}
          \sigma(t) \;=\; \sigma_0\,\sqrt{1 + \left(\frac{t}{2\sigma_0^2}\right)^2},
          \label{eq:spreading}
        \end{equation}
        che per $t \gg 2\sigma_0^2 = 0{,}5$\,u.a.\ diventa lineare:
        $\sigma(t) \approx t/(2\sigma_0)$.

  \item \textbf{Dopo lo scattering}: La larghezza cresce più velocemente
        del risultato in spazio libero. La ragione fisica è che la funzione
        d'onda si divide in due componenti spazialmente separate (riflessa
        e trasmessa), e la varianza $\langle x^2\rangle - \langle x\rangle^2$
        riceve un grande contributo dalla \emph{distanza tra i due picchi},
        non solo dalla larghezza di ciascun singolo picco. Questo è un
        effetto puramente quantistico di interferenza.
\end{itemize}

\paragraph{Cosa osservare.}
Un netto cambiamento di pendenza al momento dello scattering: il tasso di
crescita di $\sigma(t)$ aumenta bruscamente quando il pacchetto si divide.

\subsection{Animazione (\texttt{evolution.mp4})}

L'animazione mostra simultaneamente $|\psi(x,t)|^2$ (linea continua blu) e
$\mathrm{Re}[\psi(x,t)]$ (tratteggio arancione), che rivela la struttura di
fase oscillatoria della funzione d'onda.

\paragraph{Comportamento atteso.}
\begin{itemize}
  \item La parte reale oscilla alla frequenza spaziale $k_0$, visibile come
        rapide oscillazioni all'interno dell'inviluppo del pacchetto.
  \item Quando il pacchetto si avvicina alla barriera, le oscillazioni nella
        regione della barriera vengono soppresse esponenzialmente (onda evanescente).
  \item Dopo l'interazione, il pacchetto riflesso mostra oscillazioni che
        viaggiano nella direzione $-x$ (fase decrescente nel tempo), mentre
        il piccolo pacchetto trasmesso oscilla nella direzione $+x$.
  \item L'interferenza tra le componenti incidente e riflessa è visibile
        come uno schema di onda stazionaria a sinistra della barriera durante
        la fase di scattering.
\end{itemize}

% ═══════════════════════════════════════════════════════════════════════════
\section{Parametri numerici: stabilità e accuratezza}
% ═══════════════════════════════════════════════════════════════════════════

\subsection{Requisiti di risoluzione della griglia}

Affinché la simulazione rappresenti fedelmente il pacchetto d'onde, la
griglia deve risolvere le scale di lunghezza fisiche più corte:

\begin{itemize}
  \item \textbf{Spaziale}: la lunghezza d'onda di de Broglie $\lambda = 2\pi/k_0$
        deve essere risolta con diversi punti di griglia. Per $k_0 = 5$,
        $\lambda \approx 1{,}26$\,bohr, e $\dx = 0{,}01$\,bohr fornisce
        $\approx 126$ punti per lunghezza d'onda — più che sufficiente.

  \item \textbf{All'interno della barriera}: la lunghezza di decadimento
        evanescente è $1/\kappa \approx 0{,}26$\,bohr per i parametri di
        default. Con $\dx = 0{,}01$, ci sono $\approx 26$ punti per lunghezza
        di decadimento, il che è adeguato.

  \item \textbf{Temporale}: lo schema \CN{} è incondizionatamente stabile,
        quindi $\dt$ è vincolato solo dall'\emph{accuratezza}. Un criterio
        pratico è $\dt \ll 1/E_\text{max}$, dove
        $E_\text{max} \approx k_\text{max}^2/2$ è l'energia massima dello
        spettro. Per $k_\text{max} \approx k_0 + 3/(2\sigma)
        \approx 8$, $E_\text{max} \approx 32$\,u.a., ottenendo
        $\dt \ll 0{,}03$\,u.a. Il valore di default $\dt = 0{,}001$\,u.a.\ soddisfa
        facilmente questo criterio.
\end{itemize}

\subsection{Numero di Courant}

Sebbene lo schema \CN{} sia incondizionatamente stabile, il numero di
Courant effettivo $\nu = k_0\dt/\dx = 5 \times 0{,}001 / 0{,}01 = 0{,}5$
fornisce un'indicazione dell'accuratezza della velocità di fase. Valori
$\nu \lesssim 1$ garantiscono che la velocità di gruppo numerica sia vicina
a quella analitica.

% ═══════════════════════════════════════════════════════════════════════════
\section{Conclusioni}
% ═══════════════════════════════════════════════════════════════════════════

Lo schema \CN{} fornisce un metodo robusto, del secondo ordine e
conservativo della norma per propagare la TDSE. La sua natura implicita
elimina la restrizione sul passo temporale che affligge gli schemi espliciti,
e la struttura tridiagonale del sistema lineare risultante rende ogni
passo $\mathcal{O}(N)$ tramite l'algoritmo di Thomas.

La simulazione riproduce le caratteristiche quantistiche fondamentali dello
scattering da barriera: allargamento del pacchetto d'onde, effetto tunnel
attraverso una regione classicamente proibita, e la divisione del pacchetto
in componenti riflessa e trasmessa. Il confronto quantitativo del coefficiente
di trasmissione calcolato numericamente con la formula analitica a onda piana
fornisce una validazione diretta sia del modello fisico che dell'implementazione
numerica.

\clearpage

% ═══════════════════════════════════════════════════════════════════════════
\begin{thebibliography}{99}

\bibitem{Griffiths2018}
D.~J.~Griffiths e D.~F.~Schroeter,
\textit{Introduction to Quantum Mechanics}, 3ª ed.\
Cambridge University Press, 2018.

\bibitem{CrankNicolson1947}
J.~Crank e P.~Nicolson,
``A practical method for numerical evaluation of solutions of partial
differential equations of the heat-conduction type,''
\textit{Proc.\ Camb.\ Phil.\ Soc.}, vol.~43, pp.~50--67, 1947.

\bibitem{Thomas1949}
L.~H.~Thomas,
``Elliptic problems in linear difference equations over a network,''
Watson Sci.\ Comput.\ Lab.\ Report, Columbia University, 1949.

\bibitem{Tannor2007}
D.~J.~Tannor,
\textit{Introduction to Quantum Mechanics: A Time-Dependent Perspective},
University Science Books, 2007.

\bibitem{NumericalRecipes}
W.~H.~Press, S.~A.~Teukolsky, W.~T.~Vetterling e B.~P.~Flannery,
\textit{Numerical Recipes: The Art of Scientific Computing}, 3ª ed.\
Cambridge University Press, 2007.

\bibitem{Sakurai2017}
J.~J.~Sakurai e J.~Napolitano,
\textit{Modern Quantum Mechanics}, 3ª ed.\
Cambridge University Press, 2021.

\bibitem{Cohen1977}
C.~Cohen-Tannoudji, B.~Diu e F.~Lalo\"{e},
\textit{Quantum Mechanics}, voll.~I e II,
Wiley, 1977.

\bibitem{Hairer1993}
E.~Hairer, S.~P.~N\o rsett e G.~Wanner,
\textit{Solving Ordinary Differential Equations I: Nonstiff Problems},
2ª ed., Springer, 1993.

\bibitem{Leforestier1991}
C.~Leforestier et al.,
``A comparison of different propagation schemes for the time-dependent
Schr\"{o}dinger equation,''
\textit{J.\ Comput.\ Phys.}, vol.~94, pp.~59--80, 1991.

\bibitem{Wigner1955}
E.~P.~Wigner,
``Lower limit for the energy derivative of the scattering phase shift,''
\textit{Phys.\ Rev.}, vol.~98, pp.~145--147, 1955.

\end{thebibliography}

\clearpage

% ═══════════════════════════════════════════════════════════════════════════
\appendix
% ═══════════════════════════════════════════════════════════════════════════

% ───────────────────────────────────────────────────────────────────────────
\section{Lo spazio degli stati quantistici e il formalismo di Dirac}
\label{app:hilbert}
% ───────────────────────────────────────────────────────────────────────────

\subsection{Struttura dello spazio di Hilbert}

Il contesto matematico rigoroso della meccanica quantistica è uno
\emph{spazio di Hilbert complesso separabile} $\mathcal{H}$. Per una
particella priva di spin in una dimensione, la scelta naturale è
$\mathcal{H} = L^2(\mathbb{R})$: lo spazio delle funzioni complesse
a quadrato integrabile su $\mathbb{R}$, dotato del prodotto interno
\begin{equation}
  \langle f, g\rangle \;=\; \int_{-\infty}^{+\infty} f^*(x)\,g(x)\,dx.
\end{equation}
Uno stato fisico corrisponde a un \emph{vettore unitario} (raggio) in
$\mathcal{H}$, cioè $\|\psi\| = 1$.

Nella notazione di Dirac, gli stati sono \emph{ket} $\ket{\psi} \in \mathcal{H}$
e i loro duali sono \emph{bra} $\bra{\psi} \in \mathcal{H}^*$.
Il prodotto interno è $\braket{\phi}{\psi}$, e la norma è
$\|\psi\|^2 = \braket{\psi}{\psi}$.

\subsection{Osservabili come operatori autoaggiunti}

Ogni osservabile fisica $A$ corrisponde a un operatore \emph{autoaggiunto}
(hermitiano) $\hat{A}: \mathcal{H} \to \mathcal{H}$, definito da
$\hat{A}^\dagger = \hat{A}$. L'autoaggiuntezza garantisce:
\begin{itemize}
  \item Tutti gli autovalori sono reali (le quantità misurabili devono essere reali).
  \item Gli autovettori corrispondenti ad autovalori distinti sono ortogonali.
  \item Gli autovettori formano una base ortonormale completa di $\mathcal{H}$
        (teorema spettrale).
\end{itemize}

Per spettri continui (posizione, momento), gli autostati sono distribuzioni
piuttosto che vere funzioni $L^2$. Gli autostati di posizione $\ket{x'}$
soddisfano
\begin{equation}
  \hat{x}\ket{x'} = x'\ket{x'}, \qquad
  \braket{x}{x'} = \delta(x-x'), \qquad
  \int_{-\infty}^{+\infty} \ket{x'}\bra{x'}\,dx' = \hat{I},
\end{equation}
e la funzione d'onda è la proiezione
$\psi(x,t) = \braket{x}{\psi(t)}$.

\subsection{La relazione di commutazione canonica}

Gli operatori di posizione e momento soddisfano
\begin{equation}
  [\hat{x},\, \hat{p}] \;=\; \ii\hbar\,\hat{I}.
  \label{eq:ccr}
\end{equation}
Nella rappresentazione delle posizioni, $\hat{p} = -\ii\hbar\,\partial_x$.
La relazione di commutazione \eqref{eq:ccr} è la radice algebrica del
principio di indeterminazione di Heisenberg: per qualsiasi stato $\ket{\psi}$,
\begin{equation}
  \Delta x\,\Delta p \;\geq\; \frac{\hbar}{2},
\end{equation}
dove $(\Delta A)^2 = \langle \hat{A}^2\rangle - \langle\hat{A}\rangle^2$.
Il pacchetto d'onde gaussiano \eqref{eq:gauss} satura questo limite,
rendendolo lo \emph{stato coerente} dell'algebra di Weyl--Heisenberg.

\subsection{Operatore di evoluzione temporale e unitarietà}

La soluzione formale della TDSE è
\begin{equation}
  \ket{\psi(t)} \;=\; \hat{U}(t,t_0)\,\ket{\psi(t_0)},
  \qquad
  \hat{U}(t,t_0) = \exp\!\left[-\frac{\ii}{\hbar}\hat{H}(t-t_0)\right],
\end{equation}
dove l'ultima espressione vale per $\hat{H}$ indipendente dal tempo.
L'operatore $\hat{U}$ è \emph{unitario}: $\hat{U}^\dagger \hat{U} = \hat{I}$,
che è l'enunciato operatoriale della conservazione della probabilità.
Il propagatore \CN{} $A^{-1}B$ è un'approssimazione unitaria finito-dimensionale
di $\hat{U}(\dt)$.

% ───────────────────────────────────────────────────────────────────────────
\section{Metodo della matrice di trasferimento e formule esatte per il tunnel}
\label{app:transfer}
% ───────────────────────────────────────────────────────────────────────────

\subsection{Soluzione a tratti dell'equazione di Schr\"{o}dinger stazionaria}

Per un potenziale statico la TDSE si separa mediante $\psi(x,t) =
\phi(x)\,\ee^{-\ii Et/\hbar}$, dando luogo all'equazione di Schr\"{o}dinger
\emph{indipendente dal tempo} (TISE):
\begin{equation}
  -\frac{1}{2}\phi''(x) + V(x)\phi(x) \;=\; E\,\phi(x).
\end{equation}

Per la barriera rettangolare \eqref{eq:barrier}, la TISE ha coefficienti
costanti a tratti. In ciascuna regione la soluzione generale è:
\begin{align}
  \text{Regione I } (x < x_\text{bar}):&\quad
    \phi(x) = A\,\ee^{\ii kx} + B\,\ee^{-\ii kx}, \quad k = \sqrt{2E}, \\
  \text{Regione II } (x_\text{bar} < x < x_\text{bar}+w):&\quad
    \phi(x) = C\,\ee^{\kappa x} + D\,\ee^{-\kappa x}, \quad
    \kappa = \sqrt{2(V_0-E)},\quad (E<V_0), \\
  \text{Regione III } (x > x_\text{bar}+w):&\quad
    \phi(x) = F\,\ee^{\ii kx},
\end{align}
dove $F\ee^{\ii kx}$ è l'onda trasmessa (nessuna componente che si muove
verso sinistra nella regione III per un'onda incidente da destra).

\subsection{La matrice di trasferimento}

A ciascuna interfaccia, la continuità di $\phi$ e $\phi'$ fornisce due
condizioni di raccordo. Queste possono essere codificate in modo compatto
in una \emph{matrice di trasferimento} $M$ $2\times 2$ tale che
\begin{equation}
  \begin{pmatrix} A \\ B \end{pmatrix}
  \;=\; M \begin{pmatrix} F \\ 0 \end{pmatrix}.
\end{equation}
Per la barriera rettangolare, $M = M_{\text{I}\to\text{II}}\cdot
M_{\text{II}\to\text{III}}$, dove ciascuna matrice di interfaccia è
\begin{equation}
  M_{j\to k} \;=\; \frac{1}{2q_k}
  \begin{pmatrix}
    q_k + q_j & q_k - q_j \\
    q_k - q_j & q_k + q_j
  \end{pmatrix},
\end{equation}
con $q_j$ il vettore d'onda (eventualmente immaginario) nella regione $j$.
L'ampiezza di trasmissione è $t = 1/M_{11}$ e
\begin{equation}
  T = |t|^2 = \frac{1}{|M_{11}|^2},
\end{equation}
che, dopo calcolo esplicito, riproduce \eqref{eq:T_tunnelling}.
Il formalismo della matrice di trasferimento si generalizza direttamente
a potenziali costanti a tratti arbitrari.

\subsection{Approssimazione WKB}

Per una barriera di potenziale \emph{liscia}, la matrice di trasferimento
esatta non è disponibile in forma chiusa. L'approssimazione di
Wentzel--Kramers--Brillouin (WKB) sostituisce la soluzione evanescente
esatta all'interno della barriera con un ansatz localmente a onda piana:
\begin{equation}
  \phi(x) \;\approx\;
  \frac{C}{\sqrt{\kappa(x)}}\,\exp\!\left[-\int_{x_1}^{x}\kappa(x')\,dx'\right],
  \quad \kappa(x) = \sqrt{2(V(x)-E)},
\end{equation}
dove $x_1$, $x_2$ sono i punti di inversione classici ($V(x_{1,2}) = E$).
Il coefficiente di trasmissione WKB è
\begin{equation}
  T_\text{WKB} \;\approx\; \exp\!\left[-2\int_{x_1}^{x_2}\kappa(x)\,dx\right].
  \label{eq:wkb}
\end{equation}
Per la barriera rettangolare, $\kappa(x) = \kappa = \text{cost}$ e
\eqref{eq:wkb} dà $T \approx \ee^{-2\kappa w}$, coerente con il limite
di $\kappa w$ grande di \eqref{eq:T_tunnelling}.

% ───────────────────────────────────────────────────────────────────────────
\section{Relazione di dispersione e dinamica del pacchetto libero}
\label{app:dispersion}
% ───────────────────────────────────────────────────────────────────────────

\subsection{Relazione di dispersione}

Per una particella libera ($V=0$), sostituendo un'onda piana
$\psi = \ee^{\ii(kx - \omega t)}$ nella TDSE si ottiene la
\emph{relazione di dispersione}
\begin{equation}
  \omega(k) \;=\; \frac{k^2}{2} \quad \text{(unità atomiche)}.
\end{equation}
Questa è una dispersione \emph{quadratica}, a differenza della dispersione
lineare dei fotoni. Due velocità caratterizzano la propagazione del pacchetto
d'onde:
\begin{align}
  v_\text{fase} &= \frac{\omega}{k} = \frac{k}{2}, \\
  v_\text{gruppo} &= \frac{d\omega}{dk} = k.
\end{align}
La velocità di gruppo $v_g = k_0$ è la velocità del centro di massa del
pacchetto e coincide con la velocità classica $p/m = k_0$ in unità atomiche.

\subsection{Propagazione esatta del pacchetto gaussiano nello spazio libero}

Un pacchetto d'onde gaussiano nello spazio libero può essere propagato
esattamente usando il teorema di convoluzione di Fourier. Partendo da
\eqref{eq:gauss}, la sua trasformata di Fourier è
\begin{equation}
  \tilde{\phi}(k,0) \;=\;
  (2\pi\sigma^2)^{1/4}\,\sigma\,\ee^{-(k-k_0)^2\sigma^2}\,\ee^{-\ii k x_0}.
\end{equation}
Al tempo $t$, ciascun modo acquista una fase $\ee^{-\ii\omega(k)t}$:
\begin{equation}
  \tilde{\phi}(k,t) \;=\; \tilde{\phi}(k,0)\,\ee^{-\ii k^2 t/2}.
\end{equation}
Antitrasformando si ottiene
\begin{equation}
  \psi(x,t) \;=\;
  \left(\frac{1}{2\pi\sigma_t^2}\right)^{1/4}
  \exp\!\left[-\frac{(x - x_0 - k_0 t)^2}{4\sigma_t^2}\right]
  \exp\!\left[\ii k_0\!\left(x - \frac{k_0 t}{2}\right)\right]
  \exp\!\left[\ii\phi_t\right],
  \label{eq:free_packet}
\end{equation}
dove la \emph{larghezza complessa} è
\begin{equation}
  \sigma_t \;=\; \sigma\sqrt{1 + \ii t/(2\sigma^2)},
\end{equation}
e $\phi_t = -\frac{1}{2}\arg(\sigma_t/\sigma)$ è una fase globale.
La densità di probabilità è ancora gaussiana:
\begin{equation}
  |\psi(x,t)|^2 \;=\;
  \frac{1}{\sqrt{2\pi}|\sigma_t|}
  \exp\!\left[-\frac{(x-x_0-k_0 t)^2}{2|\sigma_t|^2}\right],
\end{equation}
con larghezza reale $|\sigma_t| = \sigma\sqrt{1 + t^2/(4\sigma^4)}$,
che è esattamente l'equazione \eqref{eq:spreading}.

Il tasso di allargamento $d|\sigma_t|/dt \to 1/(2\sigma)$ per $t$ grande
è inversamente proporzionale alla larghezza iniziale: un pacchetto stretto
(piccolo $\sigma$, grande dispersione di momento) si allarga più velocemente.

% ───────────────────────────────────────────────────────────────────────────
\section{Splitting degli operatori e la forma di Cayley}
\label{app:splitting}
% ───────────────────────────────────────────────────────────────────────────

\subsection{Il propagatore esatto come esponenziale di matrice}

La soluzione formale della TDSE su un passo temporale è
\begin{equation}
  \boldsymbol{\psi}^{n+1} \;=\;
  \ee^{-\ii\dt H_\text{disc}}\,\boldsymbol{\psi}^n.
\end{equation}
Il calcolo diretto dell'esponenziale di matrice costa $\mathcal{O}(N^3)$
— proibitivo per $N$ grande. Esistono diverse approssimazioni efficienti.

\subsection{Approssimanti di Padé e la forma di Cayley}

Lo schema \CN{} è equivalente all'\emph{approssimante di Padé del primo ordine}
dell'esponenziale di matrice:
\begin{equation}
  \ee^{-\ii\dt H} \;\approx\;
  \frac{I - \ii\dt H/2}{I + \ii\dt H/2}
  \;=\; A^{-1}B.
  \label{eq:cayley}
\end{equation}
Questo è noto anche come \emph{forma di Cayley} del propagatore. È
\emph{esattamente unitario} indipendentemente da $\dt$, perché il numeratore
e il denominatore sono coniugati complessi l'uno dell'altro (per $H$ hermitiano).

Approssimanti di Padé di ordine superiore $[m/m]$ esistono e raggiungono
un'accuratezza $\mathcal{O}(\dt^{2m})$ pur rimanendo unitari, ma richiedono
la soluzione di $m$ sistemi tridiagonali per passo e sono raramente usati
in pratica.

\subsection{Splitting di Strang}

Per Hamiltoniani della forma $H = T + V$ dove $T$ è l'operatore cinetico
e $V$ è diagonale nello spazio delle posizioni, il metodo di
\emph{split-operator} (splitting di Strang) approssima:
\begin{equation}
  \ee^{-\ii\dt(T+V)} \;\approx\;
  \ee^{-\ii\dt V/2}\,\ee^{-\ii\dt T}\,\ee^{-\ii\dt V/2}
  + \mathcal{O}(\dt^3).
\end{equation}
Ciascun fattore può essere applicato in modo efficiente: $\ee^{-\ii\dt V/2}$
è una moltiplicazione puntuale nello spazio delle posizioni (diagonale),
mentre $\ee^{-\ii\dt T}$ è una moltiplicazione nello spazio del momento
(di Fourier). Usando la Trasformata di Fourier Veloce (FFT), ogni passo
costa $\mathcal{O}(N\log N)$, rispetto a $\mathcal{O}(N)$ per l'algoritmo
di Thomas \CN{}. Tuttavia, il metodo split-operator è del secondo ordine
e unitario esatto solo nel limite dell'aritmetica esatta, e richiede FFT
globali che sono più difficili da parallelizzare.

% ───────────────────────────────────────────────────────────────────────────
\section{Teoria dello scattering: matrice S e tempo di ritardo di Wigner}
\label{app:scattering}
% ───────────────────────────────────────────────────────────────────────────

\subsection{La matrice di scattering}

Per un potenziale che si annulla per $x \to \pm\infty$, la teoria dello
scattering descrive la relazione asintotica tra onde piane entranti e
uscenti. In 1D, la \emph{matrice S} è una matrice unitaria $2\times 2$:
\begin{equation}
  S \;=\;
  \begin{pmatrix}
    t   & r' \\
    r   & t'
  \end{pmatrix},
\end{equation}
dove $t, t'$ sono le ampiezze di trasmissione (incidente da destra e da
sinistra) e $r, r'$ sono le ampiezze di riflessione. L'unitarietà
$S^\dagger S = I$ implica $|t|^2 + |r|^2 = 1$, cioè $T + R = 1$.

Per un potenziale simmetrico (con $V(x)$ pari), $t = t'$ e $|r| = |r'|$.
La barriera rettangolare è simmetrica rispetto al suo centro, quindi
entrambe le ampiezze di trasmissione sono uguali.

\subsection{Sfasamenti e ritardi temporali}

L'ampiezza di trasmissione può essere scritta come
$t = \sqrt{T}\,\ee^{\ii\varphi_t(k)}$, dove $\varphi_t$ è la fase di
trasmissione. Il \emph{tempo di ritardo di Wigner--Smith} è definito come
\begin{equation}
  \tau(k) \;=\; \frac{d\varphi_t}{d\omega}
  \;=\; \frac{1}{v_g}\frac{d\varphi_t}{dk},
  \label{eq:wigner}
\end{equation}
e rappresenta il tempo extra trascorso dal pacchetto d'onde trasmesso
vicino alla barriera rispetto alla propagazione libera \cite{Wigner1955}.
Per la barriera rettangolare nel regime di tunnel, $\tau > 0$ (il pacchetto
è ritardato), e per $V_0 \to \infty$ il ritardo si satura a
$\tau_\text{max} = w/v_g$ — il cosiddetto \emph{effetto Hartman}, in cui
il ritardo diventa indipendente dalla larghezza della barriera per barriere
opache.

Nella simulazione, il ritardo si manifesta come una riduzione della pendenza
di $\langle x\rangle(t)$ durante l'evento di scattering (Pannello C di
\texttt{analysis.png}).

\subsection{Il teorema ottico in 1D}

Il teorema ottico in 1D mette in relazione l'ampiezza di scattering in
avanti con la sezione d'urto totale. In una dimensione, la ``sezione
d'urto'' è sostituita dalla \emph{riflettanza totale}, e il teorema afferma:
\begin{equation}
  \mathrm{Im}(r) + \frac{k}{2}(T-1) = 0,
\end{equation}
che è una diretta conseguenza dell'unitarietà della matrice $S$ e fornisce
un ulteriore controllo di consistenza per i calcoli numerici di scattering.

% ───────────────────────────────────────────────────────────────────────────
\section{Analisi di stabilità di von Neumann}
\label{app:stability}
% ───────────────────────────────────────────────────────────────────────────

\subsection{Analisi di von Neumann per la TDSE}

La stabilità degli schemi alle differenze finite per PDE lineari può essere
analizzata rigorosamente tramite il metodo di von Neumann (di Fourier).
Si consideri una griglia uniforme e un singolo modo di Fourier
$\psi_j^n = g^n\,\ee^{\ii j\xi\dx}$, dove $\xi$ è il vettore d'onda e
$g$ è il fattore di amplificazione. La stabilità richiede $|g| \leq 1$
per ogni $\xi$.

\paragraph{Eulero in avanti (esplicito):}
Sostituendo in $\psi_j^{n+1} = \psi_j^n - \ii\dt H_\text{disc}\psi_j^n$
si ottiene
\begin{equation}
  g \;=\; 1 - \ii\dt\,\lambda(\xi),
  \quad \lambda(\xi) = \frac{2(1-\cos(\xi\dx))}{\dx^2} + V \in \mathbb{R}.
\end{equation}
Allora $|g|^2 = 1 + \dt^2\lambda^2 > 1$ per $\lambda \neq 0$: lo schema
è \emph{sempre instabile} per l'equazione di Schrödinger.

\paragraph{Crank--Nicolson:}
Il fattore di amplificazione è
\begin{equation}
  g \;=\; \frac{1 - \ii\dt\lambda/2}{1 + \ii\dt\lambda/2}.
\end{equation}
Poiché $\lambda \in \mathbb{R}$, si ha $|g|^2 = 1$ identicamente per
ogni $\xi$ e per ogni $\dt > 0$: lo schema è \emph{incondizionatamente
stabile} e, inoltre, \emph{esattamente conservativo della norma} al
livello del fattore di amplificazione.

\paragraph{Regola del punto medio implicita (Crank--Nicolson è un caso speciale):}
Più in generale, il metodo $\theta$
\begin{equation}
  \ii\frac{\psi^{n+1}-\psi^n}{\dt} \;=\;
  (1-\theta)H\psi^n + \theta H\psi^{n+1}
\end{equation}
è stabile per $\theta \geq 1/2$. Lo schema di Crank--Nicolson
corrisponde a $\theta = 1/2$ ed è l'unica scelta in questa famiglia
che sia sia stabile che del secondo ordine.

\subsection{Dispersione numerica}

Sebbene lo schema \CN{} sia incondizionatamente stabile, introduce
un \emph{errore di dispersione numerica}: la velocità di fase discreta
$\tilde{\omega}(\xi)$ differisce da quella esatta $\omega(\xi) = \xi^2/2$.
Dal fattore di amplificazione,
\begin{equation}
  \tilde{\omega}(\xi) \;=\;
  \frac{2}{\dt}\arctan\!\left(\frac{\dt\,\lambda(\xi)}{2}\right),
  \quad
  \lambda(\xi) = \frac{2(1-\cos\xi\dx)}{\dx^2}.
\end{equation}
Espandendo per piccoli $\xi\dx$ e $\dt$:
\begin{equation}
  \tilde{\omega}(\xi) \;=\; \frac{\xi^2}{2}
  \left[1 - \frac{(\xi\dx)^2}{12} - \frac{\dt^2\xi^4}{24} + \ldots\right],
\end{equation}
mostrando errori di dispersione $\mathcal{O}(\dx^2, \dt^2)$, coerenti con
l'accuratezza del secondo ordine dello schema. Le componenti ad alto $k$
del pacchetto d'onde viaggiano a una velocità leggermente errata — questa
è l'origine fisica dell'\emph{allargamento numerico} che si aggiunge
all'allargamento quantistico fisico.

% ───────────────────────────────────────────────────────────────────────────
\section{Il teorema di Ehrenfest e la corrispondenza quantistico-classica}
\label{app:ehrenfest}
% ───────────────────────────────────────────────────────────────────────────

\subsection{Enunciato e dimostrazione}

Il \emph{teorema di Ehrenfest} stabilisce la connessione tra i valori
di aspettazione quantistici e le equazioni del moto classiche. Per
qualsiasi osservabile $\hat{A}$, l'equazione del moto di Heisenberg è:
\begin{equation}
  \frac{d}{dt}\langle\hat{A}\rangle \;=\;
  \frac{1}{\ii\hbar}\langle[\hat{A},\hat{H}]\rangle
  + \left\langle\pdv{\hat{A}}{t}\right\rangle.
\end{equation}
Applicando questo a $\hat{x}$ e $\hat{p}$ con $\hat{H} = \hat{p}^2/2m + V(\hat{x})$:
\begin{align}
  \frac{d\langle\hat{x}\rangle}{dt}
  &= \frac{\langle\hat{p}\rangle}{m}, \label{eq:ehr1} \\
  \frac{d\langle\hat{p}\rangle}{dt}
  &= -\left\langle\frac{\partial V}{\partial x}\right\rangle. \label{eq:ehr2}
\end{align}
Questi sono gli analoghi quantistici delle equazioni di Hamilton. Si noti
la differenza cruciale dal risultato classico: \eqref{eq:ehr2} coinvolge
$\langle V'(\hat{x})\rangle$, non $V'(\langle\hat{x}\rangle)$.
I due coincidono solo se $V$ è al più quadratico, o se il pacchetto d'onde
è così stretto che $\psi(x,t) \approx \delta(x - x_\text{cl}(t))$.

\subsection{Implicazioni per la simulazione}

Nella prima fase di propagazione, il pacchetto d'onde è stretto rispetto
alla scala su cui $V$ varia (la barriera), quindi il teorema di Ehrenfest
prevede un moto quasi classico: $\langle x\rangle(t) \approx
x_0 + k_0 t$. Questo è esattamente il regime lineare visibile nel Pannello C
prima dell'evento di scattering.

Quando il pacchetto interagisce con la barriera, $\langle V'\rangle \neq
V'(\langle x\rangle)$, l'approssimazione classica si rompe, e gli effetti
quantistici (tunnel, interferenza) dominano. Dopo lo scattering, la funzione
d'onda consiste in due pacchetti ben separati, e $\langle x\rangle$ diventa
una media pesata che non ha corrispettivo classico.

\subsection{Conservazione dell'energia}

Dal teorema di Ehrenfest, l'energia media $\langle\hat{H}\rangle$ è
conservata:
\begin{equation}
  \frac{d}{dt}\langle\hat{H}\rangle \;=\;
  \frac{1}{\ii\hbar}\langle[\hat{H},\hat{H}]\rangle = 0.
\end{equation}
Nella simulazione numerica, la conservazione dell'energia non è imposta
esplicitamente dallo schema \CN{} (solo la norma è esattamente conservata).
Tuttavia, l'accuratezza del secondo ordine dello schema garantisce che
la deriva dell'energia sia $\mathcal{O}(\dt^2)$ per passo.

% ───────────────────────────────────────────────────────────────────────────
\section{Condizioni al contorno assorbenti e perfectly matched layers}
\label{app:abc}
% ───────────────────────────────────────────────────────────────────────────

\subsection{Il problema con le condizioni al contorno di Dirichlet}

Le condizioni al contorno di Dirichlet $\psi(x_\text{min}) = \psi(x_\text{max}) = 0$
mimano pareti infinitamente rigide. Sono semplici da implementare ma introducono
\emph{riflessioni spurie} non appena il pacchetto d'onde raggiunge il bordo.
Il dominio deve quindi essere sufficientemente grande che nessuna parte di
$\psi$ raggiunga le pareti durante la simulazione — un requisito che può
rendere il problema oneroso per tempi lunghi o pacchetti veloci.

\subsection{Condizioni al contorno assorbenti}

Un'alternativa più efficiente è modificare l'Hamiltoniano vicino ai bordi
aggiungendo un potenziale immaginario negativo (assorbimento):
\begin{equation}
  V(x) \;\to\; V(x) - \ii\,\Gamma(x),
  \quad \Gamma(x) \geq 0,
\end{equation}
dove $\Gamma(x)$ è grande solo vicino ai bordi del dominio. Questo rende
l'Hamiltoniano \emph{non-hermitiano} nella regione assorbente, causando
la diminuzione della norma man mano che la probabilità viene assorbita.
L'onda numerica riflessa viene soppressa esponenzialmente.

Una scelta standard è la rampa a coseno:
\begin{equation}
  \Gamma(x) \;=\;
  \Gamma_0\,\sin^2\!\left(\frac{\pi}{2}
  \cdot\frac{|x|-x_\text{abs}}{L_\text{abs}}\right)
  \quad \text{per } |x| > x_\text{abs},
\end{equation}
dove $L_\text{abs}$ e $\Gamma_0$ sono la larghezza e la forza dello strato
di assorbimento.

\subsection{Perfectly matched layers}

Lo standard per le condizioni al contorno assorbenti è il \emph{perfectly
matched layer} (PML), originariamente sviluppato per le equazioni di Maxwell
\cite{NumericalRecipes} e successivamente adattato alla TDSE. Un PML viene
costruito prolungando analiticamente la coordinata spaziale nel piano
complesso:
\begin{equation}
  x \;\to\; \tilde{x}(x) \;=\; x + \ii\int_0^x \sigma(x')\,dx',
\end{equation}
dove $\sigma(x') \geq 0$ è il profilo di assorbimento del PML. Questa
trasformazione preserva la forma dell'equazione di Schr\"{o}dinger nella
regione PML rendendo al contempo tutte le onde uscenti evanescenti — e,
in modo cruciale, \emph{senza creare alcuna riflessione all'interfaccia PML},
indipendentemente dall'angolo o dalla frequenza dell'onda incidente.

Per la TDSE 1D, la modifica PML è equivalente all'introduzione di un
fattore di stiramento coordinato complesso $s(x)$:
\begin{equation}
  \Htop_\text{PML} \;=\;
  -\frac{1}{2s(x)}\frac{\partial}{\partial x}\frac{1}{s(x)}\frac{\partial}{\partial x}
  + V(x),
  \quad s(x) = 1 + \ii\sigma(x)/\omega.
\end{equation}
La simulazione attuale utilizza semplici condizioni al contorno di Dirichlet,
che sono sufficienti per i parametri di default. L'implementazione di
condizioni al contorno assorbenti permetterebbe simulazioni su domini molto
più piccoli o su tempi molto più lunghi.

\end{document}